
%%%%%%%%%%%%%%%%%%%%%%%%%%%%%%%%%%%%%%%%%%%%%%%%%%%%%%%
\begin{frame}{Analyse des coûts}{

Avec l'utilitaire \texttt{TKPROF}, on obtient\,:

\begin{itemize}
    \item Le coût CPU
   \item Le nombre d'entrées/sorties physiques
   \item Le nombre d'entrées/sorties en mémoire
   \item Le nombre de tris
   \item etc
\end{itemize}

Un excellent outil à utiliser\,: la
trace automatique qui donne à la fois le 
\texttt{EXPLAIN} et les coûts d'exécution.
}\end{frame}

\end{document}
%------------------------------------------------------------
\section{Exercices d'application}
%------------------------------------------------------------
%%%%%%%%%%%%%%%%%%%%%%%%%%%%%%%%%%%%%%%%%%%%%%%%%%%%%%%
\begin{frame}{L'environnement}{

Sur la machine FENRIS, se connecter sous 
le compte form$<1, 2, 3, 4, 5, 6>$ (pas de mot de passe).

\begin{itemize}
   \item Récupérer les fichiers \textit{schema.sql}
                   et \textit{CrFilms} (ou bien travailler
                   avec la base EKIP)

   \item Créer une base

   \item Sous SQLPLUS, se mettre en mode trace avec la commande\,:

\texttt{ALTER SESSION SET AUTOTRACE TRACEONLY}

\end{itemize}

Et c'est parti\,!
}\end{frame}
%%%%%%%%%%%%%%%%%%%%%%%%%%%%%%%%%%%%%%%%%%%%%%%%%%%%%%%
\begin{frame}{Optimiseur en mode règles}{

Effectuer et analyser des requêtes représentatives
 sur la base \textit{Films}.

\begin{itemize}
  \item Des sélections sur des champs indexés
        ou non indexés (faire une sélection sur le genre par exemple)

  \item Des jointures (penser à inverser l'ordre des tables
          dans le \sqlcom{FROM})

  \item Des requêtes avec \texttt{NOT IN} ou \texttt{NOT EXISTS}
\end{itemize}

}\end{frame}
%%%%%%%%%%%%%%%%%%%%%%%%%%%%%%%%%%%%%%%%%%%%%%%%%%%%%%%
\begin{frame}{Optimiseur en mode coûts}{

Maintenant, analyser les tables, les index et les distributions
      avec \texttt{ANALYSE}, et regarder les modifications
des plans d'exécution en optimisation par les coûts.
\begin{itemize}
  \item Chercher des films par le genre en mode règles
  \item Chercher des films par le genre en mode coûts
\end{itemize}
Le genre étant peu sélectif, l'optimiseur ne devrait pas
utiliser l'index dans le second cas.

}\end{frame}
%%%%%%%%%%%%%%%%%%%%%%%%%%%%%%%%%%%%%%%%%%%%%%%%%%%%%%%
\begin{frame}{Requêtes imbriquées}{

\begin{itemize}
  \item Reprendre les requêtes du polycopié
cherchant le film paru en 1958 avec James Stewart.

 \item Analyser les différentes versions 
(sans imbrication, avec
imbrication mais sans corrélation (\texttt{IN}),
avec imbrication et corrélation (\texttt{EXISTS}), 
etc). 

 \item Comparer les plans d'exécution obtenus
dans chaque cas.
\end{itemize}

}\end{frame}
%%%%%%%%%%%%%%%%%%%%%%%%%%%%%%%%%%%%%%%%%%%%%%%%%%%%%%%
\begin{frame}{Fonctions de groupe}{

\begin{itemize}
  \item  Maintenant appliquez des fonctions de groupe,
des \texttt{GROUP BY} et des \texttt{HAVING},
regardez les plans d'exécution et interprétez-les.

 \item Essayer d'appliquer l'opérateur \texttt{LIKE}
dans des sélections sur des attributs indexés.
Que constate-t-on\,? 
\end{itemize}

}\end{frame}

\end{document}

\begin{frame}{Sélection conjonctive avec deux index}{


\begin{verbatim}
   SELECT  nom FROM   Salle 
    WHERE ID-cinéma = 1098 AND capacité = 150
\end{verbatim}
Plan d'exécution\,:
\begin{verbatim}
0 SELECT STATEMENT
  1 TABLE ACCESS BY ROWID SALLE
    2 AND-EQUAL
      3 INDEX RANGE SCAN IDX-SALLE-CINEMA-ID
      4 INDEX RANGE SCAN IDX-CAPACITE
\end{verbatim}
}\end{frame}

\begin{frame}{Sélection disjonctive avec index}{


\begin{tabbing}
   \=  SELECT \= nom FROM   Salle \\
   \>  WHERE  \>  ID-cinéma = 1098 OR capacité $>$ 150
\end{tabbing}
Plan d'exécution\,:
\begin{verbatim}
0 SELECT STATEMENT
  1 CONCATENATION
    2 TABLE ACCESS BY ROWID SALLE
      3 INDEX RANGE SCAN IDX-CAPACITE
    4 TABLE ACCESS BY ROWID SALLE
      5 INDEX RANGE SCAN IDX-SALLE-CINEMA-ID
\end{verbatim}
}\end{frame}

\begin{frame}{Sélection disjonctive avec et sans index}{


\begin{verbatim}
    SELECT  nom FROM   Salle 
    WHERE   ID-cinema = 1098 OR nom = 'Salle 1'
\end{verbatim}
Plan d'exécution\,:
\begin{verbatim}
0 SELECT STATEMENT
  1 TABLE ACCESS FULL SALLE 
\end{verbatim}
}\end{frame}

\begin{frame}{Jointure avec index}{


\begin{verbatim}
    SELECT Cinéma.nom,capacité FROM cinéma, salle
    WHERE  Cinéma.ID-cinéma = salle.ID-cinéma
\end{verbatim}
Plan d'exécution\,:
\begin{verbatim}
0 SELECT STATEMENT
  1 NESTED LOOPS
    2 TABLE ACCESS FULL SALLE
    3 TABLE ACCESS BY ROWID CINEMA
      4 INDEX UNIQUE SCAN IDX-CINEMA-ID
\end{verbatim}
}\end{frame}

\begin{frame}{Jointure et sélection avec index}{


\begin{tabbing}
   \= SELECT \= Cinéma.nom,capacité FROM Cinéma, Salle\\
   \=  WHERE \=  Cinema.ID-cinéma = salle.ID-cinéma \\
   \=  AND   \= capacité $>$ 150 
\end{tabbing}
Plan d'exécution\,:
\begin{verbatim}
0 SELECT STATEMENT
  1 NESTED LOOPS
    2 TABLE ACCESS BY ROWID SALLE
      3 INDEX RANGE SCAN IDX-CAPACITE
    4 TABLE ACCESS BY ROWID CINEMA
      5 INDEX UNIQUE SCAN IDX-CINEMA-ID
\end{verbatim}
}\end{frame}

\begin{frame}{Jointure sans index}{


\begin{verbatim}
    SELECT titre
    FROM Film, Séance 
    WHERE  Film.ID-film  = Séance.ID-film
    AND     heure-début = '14H00'
\end{verbatim}
Plan d'exécution\,:
\begin{verbatim}
0 SELECT STATEMENT
  1 MERGE JOIN
    2 SORT JOIN
      3 TABLE ACCESS FULL SEANCE
    4 SORT JOIN
      5 TABLE ACCESS FULL FILM
\end{verbatim}
}\end{frame}

\begin{frame}{Différence}{


Dans quel cinéma ne peut-on voir de film après 23H\,?
\begin{tabbing}
  \=  SELECT \= Cinéma.nom \\
  \>  FROM   \> Cinéma, Salle\\
  \>  WHERE  \> Cinéma.ID-cinéma = Salle.ID-cinéma\\
  \>  AND    \> NOT EXISTS \= (SELECT * FROM séance\\
 \>         \>             \> WHERE  Salle.ID-salle = Séance.ID-salle\\
 \>         \>             \> AND    heure-fin $>$ '23H00')
\end{tabbing}
}\end{frame}

\begin{frame}{Plan d'exécution donné par EXPLAIN}
{

\begin{verbatim}
0 SELECT STATEMENT
  1 FILTER
    2 NESTED LOOPS
      3 TABLE ACCESS FULL SALLE
      4 TABLE ACCESS BY ROWID CINEMA
        5 INDEX UNIQUE SCAN IDX-CINEMA-ID
    6 TABLE ACCESS BY ROWID SEANCE
      7 INDEX RANGE SCAN IDX-SEANCE-SALLE-ID
\end{verbatim}
}\end{frame}


